\documentclass[11pt,a4paper]{article}

\usepackage[utf8]{inputenc}
\usepackage[T1]{fontenc}
\usepackage[spanish,es-tabla]{babel}
\usepackage{lmodern}
\usepackage[a4paper,margin=2.5cm]{geometry}
\usepackage{microtype}
\usepackage{enumitem}
\usepackage{parskip}
\usepackage{xcolor}
\usepackage{amsmath}
\usepackage{amssymb}
\usepackage{hyperref}

\hypersetup{
  colorlinks=true,
  linkcolor=black,
  urlcolor=blue!60!black,
  pdftitle={Planificacion inicial del trabajo},
  pdfauthor={Equipo proyecto final}
}

\setlist[itemize]{leftmargin=1.2em, itemsep=0.3em, topsep=0.3em}

\title{Planificación inicial del trabajo}
\author{Proyecto final --- Métodos numéricos}
\date{}

\begin{document}

\maketitle

\tableofcontents
\vspace{1em}

\section{Introducción}

El objetivo general de este trabajo es estudiar cómo se pueden usar las redes neuronales como herramienta numérica para aproximar modelos de pricing de opciones. Es decir, no nos interesa la red como un sistema que aprende del mercado, sino como un aproximador funcional capaz de imitar a un solver matemático ya conocido. Esta distinción puede parecer sutil pero cambia completamente lo que se está haciendo: no estamos prediciendo nada, estamos sustituyendo un cálculo costoso por una aproximación rápida y diferenciable.

En finanzas cuantitativas hay muchos modelos de pricing de opciones, desde el clásico Black-Scholes hasta modelos modernos con volatilidad estocástica, saltos o memoria. Todos ellos tienen una estructura común: dado un conjunto de inputs, que pueden ser características del contrato y parámetros del modelo, devuelven un precio o una volatilidad implícita. El problema es que conforme los modelos ganan realismo también se vuelven más caros computacionalmente. Lo que en Black-Scholes es una fórmula cerrada se convierte en Heston en una integral que hay que resolver por Fourier, y en Bates o en rough volatility en algo aún más pesado. Cuando estos modelos hay que evaluarlos miles o millones de veces, por ejemplo en calibración diaria, en cálculo de Greeks o en simulación de escenarios, el coste se vuelve un cuello de botella real.

La idea de los deep surrogates es resolver ese cuello de botella entrenando una red neuronal a partir de datos sintéticos generados por el propio modelo, de forma que la red aprenda la función de pricing. Si llamamos $f: \mathcal{X} \subset \mathbb{R}^d \to \mathbb{R}$ a la función de pricing del modelo verdadero, lo que hacemos es entrenar una red $\hat{f}_\theta$ con parámetros $\theta$ para que aproxime
\begin{equation}
\hat{f}_\theta(x) \approx f(x), \qquad x \in \mathcal{X}.
\end{equation}
Una vez entrenada, evaluarla es órdenes de magnitud más rápido que llamar al solver original.

\begin{center}
\fbox{\texttt{modelo financiero $\rightarrow$ datos sintéticos $\rightarrow$ red neuronal $\rightarrow$ evaluación rápida}}
\end{center}

La red no reemplaza la teoría financiera, simplemente aprende a imitar al solver. El objetivo no es aplicar las técnicas de deep learning a los datos de mercado ruidosos, sino utilizar estas técnicas para resolver un problema de aproximación numérica.

En Black-Scholes la fórmula cerrada ya es muy rápida y un surrogate no aporta valor, pero es ideal como caso de validación, porque conocemos la solución exacta y podemos medir errores con precisión. El objetivo del proyecto es verificar si se pueden utilizar surrogates con modelos más complejos sin perder precisión.

\section{Revisión bibliográfica}

La primera fase de nuestro trabajo ha consistido en realizar una investigación bibliográfica profunda para tener una visión global de los temas de estudio más relevantes. Hay cinco papers que hemos trabajado a fondo y que forman el núcleo de la revisión: los tres primeros eran una recomendación de los profesores como punto de partida, un cuarto, el de deep surrogates, lo encontramos nosotros buscando trabajos más recientes en la línea de aproximación numérica de modelos de pricing, y un quinto, el de differential machine learning de Huge y Savine, lo hemos incorporado como complemento al anterior debido a que el tema que trata nos parece otra línea de investigación que puede ser interesante para nuestro trabajo.

El punto de partida histórico es el trabajo de Hutchinson, Lo y Poggio de 1994~\cite{hutchinson1994}. Estos autores proponen un método no paramétrico para estimar la fórmula de pricing de un derivado mediante learning networks, en una época en la que la literatura financiera estaba dominada por enfoques paramétricos basados en Black-Scholes y argumentos de no arbitraje. La idea central es invertir la lógica habitual: en lugar de partir de un modelo estocástico para el subyacente y derivar analíticamente la fórmula del precio, dejan que la red aprenda directamente la función que mapea variables económicas observables, como el precio del subyacente normalizado por el strike y el tiempo a vencimiento, al precio de la opción, sin imponer supuestos de lognormalidad ni continuidad de trayectoria. Comparan cuatro arquitecturas, redes de funciones de base radial, perceptrones multicapa, projection pursuit regression y mínimos cuadrados ordinarios como baseline, y atacan a la vez los problemas de pricing y delta-hedging, lo cual es interesante porque la cobertura exige que la red aproxime bien no solo el precio sino también su derivada respecto al subyacente. Validan primero con datos sintéticos generados por Monte Carlo bajo Black-Scholes y después con datos reales de opciones sobre futuros del S\&P 500 entre 1987 y 1991, donde la red llega a superar a Black-Scholes en tracking error de la cartera replicante. Su relevancia histórica reside en que es el primer trabajo que demuestra de forma sistemática que una red neuronal puede aprender una fórmula de pricing y usarse operativamente para cubrir, abriendo la línea de investigación que conecta directamente con el enfoque actual de deep surrogates. La diferencia clave es que aquí la red aprende del mercado, mientras que un surrogate moderno aprende la salida de un pricer ya conocido para acelerarlo, lo que cambia el papel de la red de competidor de Black-Scholes a aproximador eficiente del mismo.

El segundo paper es el de Becker, Cheridito y Jentzen sobre deep optimal stopping, de 2019~\cite{becker2019}. Aquí los autores proponen un método de deep learning para resolver problemas de optimal stopping en tiempo discreto, es decir, problemas en los que hay que decidir en cada instante si parar o continuar un proceso de Markov para maximizar la esperanza de un pago. La idea central es que cualquier tiempo de parada admisible se puede descomponer en una sucesión de decisiones binarias, una por cada paso temporal, y que cada una de esas decisiones puede aprenderse como una función del estado actual mediante una red neuronal feedforward. Sustituyen la indicadora dura por una sigmoide, lo que permite entrenar los parámetros con ascenso de gradiente estocástico sobre trayectorias Monte Carlo, y proceden por inducción hacia atrás desde el último instante hasta el primero. El método se aplica al pricing de opciones bermudas tipo max-call sobre cestas de hasta quinientos activos en un Black-Scholes multidimensional, a un callable multi barrier reverse convertible y al problema no markoviano de parar óptimamente un movimiento browniano fraccionario, con precisión alta y tiempos cortos en dimensiones donde los métodos tradicionales sufren la maldición de la dimensionalidad. Para nuestra revisión es importante porque ilustra una línea claramente distinta de la nuestra: ellos aprenden directamente la política de ejercicio para derivados con derecho de parada anticipada, mientras que en nuestro proyecto la red actuará como surrogate de la función de pricing de opciones europeas, aproximando el mapa de parámetros del modelo al precio sin necesidad de modelar decisiones de ejercicio. Son enfoques complementarios dentro del uso del deep learning en finanzas cuantitativas, pero no resuelven el mismo problema.

El tercer paper recomendado es la revisión bibliográfica de Ruf y Wang de 2020~\cite{ruf2020}, que funciona como mapa general del campo. Los autores compilan y comparan más de ciento cincuenta artículos en una tabla detallada que cataloga cada estudio según seis dimensiones, entre ellas las variables de entrada, las salidas de la red, los modelos benchmark, las métricas de rendimiento, el método de partición de los datos y los activos subyacentes. Cubren tanto el enfoque puramente no paramétrico, donde la red aprende directamente el precio a partir de variables como subyacente, strike, vencimiento y volatilidad, como los enfoques híbridos que combinan fórmulas paramétricas tipo Black-Scholes con correcciones aprendidas, además de variantes más recientes orientadas a calibración, resolución de PDEs, aproximación de funciones de valor en problemas de control óptimo y deep hedging. Lo más interesante de este paper son las lecciones que extrae del repaso exhaustivo de la bibliografía hasta la fecha: la importancia de usar inputs estacionarios como la moneyness en lugar de precio y strike por separado, la necesidad de elegir benchmarks que no trivialicen la comparación, y sobre todo el problema del data leakage que aparece cuando la partición train/test se hace de manera aleatoria sobre datos con estructura temporal, rompiendo el orden cronológico, contaminando el test e infraestimando sistemáticamente el error de generalización. Para nuestro trabajo es una referencia muy útil porque ofrece un panorama de la disciplina, permite situar nuestra propuesta dentro de las diferentes líneas de investigación históricas y nos advierte de errores típicos que conviene evitar. No obstante, la mayor parte de la bibliografía analizada en este artículo versa sobre el uso de técnicas de machine learning para estimar la función de pricing real directamente sobre datos de mercado.

El cuarto paper, y el que nos ha terminado convenciendo como referencia central, es el de Chen, Didisheim y Scheidegger de 2025 sobre deep surrogates para finanzas~\cite{chen2025}. Como ya hemos descrito en detalle en las secciones anteriores, este artículo propone los deep surrogates como una clase de aproximadores numéricos para modelos estructurales costosos, sustituyendo la función original por una red neuronal profunda entrenada con muestras simuladas que toma exactamente los mismos inputs y devuelve la misma salida a un coste computacional drásticamente menor, además de proporcionar gradientes a un coste que no escala con el número de inputs. Una idea crucial de este paper descansa en que la maldición de la dimensionalidad queda aquí mitigada en el caso de los surrogates por dos factores que se refuerzan entre sí. Por un lado, los datos de entrenamiento se generan de manera sintética y determinista a partir del propio modelo, sin ruido estadístico, lo que elimina la necesidad de cubrir el espacio de inputs con una densidad de muestras que permita promediar errores aleatorios; por otro lado, la función de pricing en modelos afines como Heston o Bates es analíticamente suave, lo que controla el problema de aproximación y permite a la red profunda representar el mapa de inputs a precio con un número de parámetros que escala de forma manejable con la dimensión. En el paper emplean una versión del modelo de Bates, en un espacio aumentado de dimensión trece, mapean los precios a Black-Scholes implied volatility para que los errores sean comparables entre opciones de distinta moneyness y vencimiento, eligen Swish frente a ReLU para que las Deltas sean estables, validan por bins de moneyness y madurez, y demuestran que el surrogate reduce la calibración diaria del modelo de las ciento veinticinco horas que necesita la inversión FFT a poco más de dos horas en CPU y unas decenas de minutos en GPU. Esa reducción de tiempos es fundamental y abre la puerta a aplicaciones empíricas que antes eran inviables, como el índice diario TailRisk propuesto para predecir crashes utilizando la información del mercado de opciones o la medida de inestabilidad de parámetros vinculada a la liquidez de mercado y competitividad de las quotes de los market makers.

El quinto paper que hemos considerado de gran interés es el de Huge y Savine de 2020 sobre differential machine learning~\cite{huge2020}. Originalmente lo teníamos como uno más del panorama, pero al pensar en cómo podríamos innovar respecto al paper de deep surrogates nos hemos dado cuenta de que ofrece una pieza metodológica muy potente y muy alineada con nuestras preocupaciones. La idea es que cuando entrenas un surrogate con datos sintéticos, normalmente la red solo ve el valor de la función en cada punto, por ejemplo el precio de la opción, y aprende a predecirlo. Sin embargo, si el modelo verdadero es diferenciable, también conoces la pendiente de la función en cada punto, es decir, las griegas. Differential machine learning aprovecha esa información añadiendo los gradientes verdaderos al conjunto de entrenamiento y modificando la pérdida para que la red no solo estime el precio, sino también su derivada respecto a los inputs. El resultado es que se consiguen aproximaciones mucho más precisas, especialmente en términos de Greeks, con datasets sustancialmente más pequeños. Para nuestro trabajo es una referencia clave porque conecta directamente con la necesidad de validar Greeks que ya teníamos heredada del paper de Deep Surrogates y porque tiene una motivación matemática desde el punto de vista de aproximación de funciones.

Más allá de estos cinco pilares, hemos echado un vistazo a otros papers que ayudan a entender que la aplicación de redes neuronales a derivados ha avanzado en varias direcciones complementarias. Una primera línea ataca de frente el problema de resolver las ecuaciones diferenciales parciales que aparecen en el pricing de derivados de alta dimensión, donde los métodos de mallado tradicionales chocan con la maldición de la dimensionalidad: en esa familia se inscribe el Deep Galerkin Method de Sirignano y Spiliopoulos de 2018~\cite{sirignano2018}, que entrena una red profunda para satisfacer el operador diferencial sobre puntos muestreados al azar y prescinde por completo de la malla, así como el Deep BSDE Solver de Han, Jentzen y E del mismo año~\cite{han2018}, que reformula la PDE como una ecuación diferencial estocástica retrógrada y aproxima el gradiente de la solución mediante redes neuronales, demostrando viabilidad incluso en problemas de cien dimensiones.

Una segunda línea busca acelerar la valoración aprendiendo directamente la función de precios, y aquí el trabajo de Ferguson y Green de 2018~\cite{ferguson2018} muestra que una red profunda entrenada con datos sintéticos generados por Monte Carlo puede valorar opciones basket millones de veces más rápido que el modelo subyacente. Una tercera vertiente aborda la cobertura de carteras como un problema de aprendizaje por refuerzo, como en el deep hedging de Buehler, Gonon, Teichmann y Wood de 2019~\cite{buehler2019}, que parametriza la estrategia de cobertura mediante redes neuronales y la optimiza directamente bajo medidas de riesgo convexas, incorporando de forma natural costes de transacción y otras fricciones de mercado sin necesidad de griegas. La calibración de modelos estocásticos también se ha beneficiado de estas técnicas, como ilustra el paper de calibración de volatilidad estocástica de 2023~\cite{dlcalib2023} al aplicar differential ML al modelo de Heston para reducir drásticamente el tiempo de calibración.

Después de revisar todo este material, la conclusión a la que hemos llegado es que el paper de deep surrogates es el que mejor encaja con nuestros objetivos. La razón principal es que el enfoque de surrogates aborda un problema que en la práctica de la industria es muy importante, mejorar el coste computacional de los modelos de pricing y reducir los tiempos de calibración y estimación de parámetros. Son temas que están muy de actualidad en mesas de derivados y en research cuantitativo, sobre todo a medida que los modelos de pricing se vuelven más realistas y por tanto más caros de evaluar. Además, desde el punto de vista de métodos numéricos el enfoque de Surrogates es el más completo e interesante.

\section{Análisis del paper de deep surrogates}

Como el paper de Chen, Didisheim y Scheidegger es la referencia central del trabajo, conviene dedicarle un análisis detallado. En las siguientes subsecciones repasamos primero qué proponen y a qué modelo lo aplican, después analizaremos por qué eligen redes neuronales frente a otras alternativas de aproximación numérica y las decisiones técnicas concretas que toman al construir el surrogate.

\subsection{¿Qué proponen en el paper?}

El paper aplica todo esto al modelo de Bates~\cite{bates1996}. Este modelo combina volatilidad estocástica al estilo Heston con saltos en el precio del activo, donde los saltos siguen una distribución doble exponencial asimétrica en lugar de la normal del Bates original. La dinámica del subyacente y la varianza bajo la medida de riesgo neutral se puede escribir como
\begin{align}
\frac{dS_t}{S_{t-}} &= (r - q - \lambda \bar{\nu})\, dt + \sqrt{v_t}\, dW^S_t + d\!\left(\sum_{i=1}^{N_t}(e^{J_i} - 1)\right), \\
dv_t &= \kappa(\theta - v_t)\, dt + \xi \sqrt{v_t}\, dW^v_t, \qquad dW^S_t\, dW^v_t = \rho\, dt,
\end{align}
donde $N_t$ es un Poisson de intensidad $\lambda$ y los tamaños de salto $J_i$ siguen una doble exponencial asimétrica
\begin{equation}
f_J(z) = p\, \eta_+ e^{-\eta_+ z}\, \mathbf{1}_{\{z\geq 0\}} + (1-p)\, \eta_- e^{\eta_- z}\, \mathbf{1}_{\{z<0\}}.
\end{equation}
Esto les permite capturar cosas que Black-Scholes no puede capturar bien, como el leverage effect, que es el fenómeno por el cual cuando el mercado cae la volatilidad suele subir, o la asimetría entre saltos positivos y negativos, algo muy importante porque en la realidad las caídas extremas y las subidas extremas no se comportan igual. La distribución doble exponencial permite modelar colas más pesadas que la normal y separar el comportamiento de saltos hacia arriba y hacia abajo.

Cuando llevan Bates al framework del surrogate, acaban con un input de trece dimensiones. Tienen cinco variables de estado, que son la moneyness normalizada, el tiempo a vencimiento, la varianza instantánea, el tipo de interés y el dividend yield, además de ocho parámetros estructurales, que controlan la velocidad de reversión a la media de la volatilidad, la varianza de largo plazo, la volatilidad de la varianza, la correlación entre precio y volatilidad, la intensidad de saltos, los tamaños medios de saltos positivos y negativos y la probabilidad de que un salto sea positivo. Formalmente, el input de la red es
\begin{equation}
x = (\,\underbrace{m, \tau, v_0, r, q}_{\text{5 estados}},\ \underbrace{\kappa, \theta, \xi, \rho, \lambda, \eta_+, \eta_-, p}_{\text{8 parámetros}}\,) \in \mathbb{R}^{13},
\end{equation}
y la red aprende una función $\hat{f}_\theta(x)$ que devuelve la volatilidad implícita Black-Scholes generada por Bates.

Aquí hay algo importante que conviene tener muy presente, y es que el tamaño final de la muestra de entrenamiento que usan es de mil millones de observaciones, $N = 10^9$. Aunque parece un número muy elevado, los autores recalcan que en un espacio de dimensión trece eso sigue siendo una muestra muy dispersa, debido a la ampliamente conocida maldición de la dimensionalidad. Esta observación es relevante porque demuestra que la red no está memorizando una tabla densa, sino que realmente está aprendiendo una estructura funcional. Aquí entra una de las justificaciones teóricas más fuertes del enfoque. Para ciertos tipos de modelos financieros, en particular los modelos afines como Heston o Bates, existen resultados teóricos que muestran que el número de parámetros entrenables y el tamaño de la muestra necesarios para aproximar bien la solución crecen polinómicamente con la dimensión del input, no exponencialmente. Esta propiedad se sustenta en el trabajo clásico de Barron de 1993 sobre la aproximación de funciones con norma de Barron acotada y, más recientemente, en los resultados de Grohs, Hornung, Jentzen y von Wurstemberger de 2018, que demuestran formalmente cómo las redes neuronales profundas escapan a la maldición de la dimensionalidad en la aproximación de soluciones de la ecuación de Black-Scholes en alta dimensión. Eso ayuda a aliviar la maldición de la dimensionalidad y es justo la razón por la que las redes neuronales son una herramienta interesante en este contexto.

\subsection{Comparación de redes neuronales con otras alternativas}

El paper compara explícitamente las redes neuronales profundas con otros métodos de aproximación que serían candidatos naturales, como polinomios, splines, sparse grids y Gaussian processes. Los polinomios funcionan razonablemente en funciones suaves pero les cuesta capturar rasgos locales fuertes como kinks, y eso en opciones es un problema serio porque el payoff tiene un kink claro en el strike. Si intentas aproximar con polinomios globales aparecen oscilaciones no deseadas, que es el famoso fenómeno de Runge. Los splines capturan bien rasgos locales porque trabajan por tramos, pero escalan muy mal en alta dimensión. Los sparse grids resuelven parcialmente el problema de las mallas cartesianas, pero funcionan peor en dominios irregulares, y los modelos financieros muchas veces tienen restricciones entre parámetros que generan dominios complicados. Los Gaussian processes son potentes y elegantes, pero escalan fatal con muchos datos porque las operaciones matriciales son del orden de $O(n^3)$, así que con millones de observaciones se vuelven inviables. Las redes neuronales cumplen los cuatro criterios de forma razonable: funcionan en alta dimensión, capturan no linealidades, admiten dominios irregulares y escalan bien con muchos datos.

Hay además otra ventaja muy importante de las redes que el paper subraya y que conecta directamente con el espíritu de métodos numéricos, que es que el Jacobiano de la red, $\partial \hat{f}_\theta / \partial x$, está disponible a un coste muy reducido mediante diferenciación automática o backpropagation. Esto es crítico en finanzas porque muchas magnitudes que importan no son el precio en sí sino sus derivadas:
\begin{equation}
\Delta = \frac{\partial C}{\partial S}, \qquad \mathcal{V} = \frac{\partial C}{\partial \sigma}, \qquad \rho = \frac{\partial C}{\partial r},
\end{equation}
todas ellas fundamentales para cobertura y gestión de riesgo. Si trabajas con el modelo original, muchas veces tienes que calcular derivadas mediante diferencias finitas,
\begin{equation}
\frac{\partial f}{\partial x_i} \approx \frac{f(x + h\, e_i) - f(x - h\, e_i)}{2h},
\end{equation}
lo que eleva el coste computacional. Con una red neuronal, tanto el cálculo de greeks como el proceso de calibración son mucho más baratos, de hecho, el paper documenta que estimar los parámetros con el método tradicional vía FFT lleva alrededor de ciento veinticinco horas, mientras que con el surrogate lleva poco más de dos horas en CPU y unas decenas de minutos en GPU. Es decir, no es una mejora marginal, es una reducción de órdenes de magnitud.

\subsection{Decisiones clave de diseño}

Hay cuatro decisiones técnicas en el paper que nos parecen especialmente importantes y que queremos incorporar a nuestro propio enfoque:

\begin{itemize}
\item \textbf{Activación suave en lugar de ReLU.} Intuitivamente uno pensaría que ReLU, definida como $\mathrm{ReLU}(x) = \max(x, 0)$, es la elección natural porque se parece al payoff de una call $(S - K)^+$, con esa forma de cero hasta cierto punto y luego crecimiento lineal. Sin embargo los autores eligen Swish, que se puede ver como una versión suave de ReLU:
\begin{equation}
\mathrm{Swish}(x) = x \cdot \mathrm{sigmoid}(\beta x) = \frac{x}{1 + e^{-\beta x}}.
\end{equation}
La razón es que ReLU no es diferenciable en el cero y, aunque eso puede no afectar mucho al precio aproximado, sí afecta a las derivadas que la red produce. En el paper se realiza un experimento muy revelador donde entrenan dos surrogates idénticos salvo por la activación y encuentran que los errores de pricing son comparables pero los errores de Delta son sustancialmente mayores con ReLU. La lección que queremos heredar es que para un surrogate financiero no basta con mirar el error en precio, también hay que mirar la calidad de las derivadas, y la suavidad de la activación importa más para los Greeks que para el precio.

\item \textbf{Volatilidad implícita como métrica de evaluación.} El paper convierte los precios a volatilidad implícita Black-Scholes y mide errores en esa escala. La volatilidad implícita $\sigma_{\mathrm{IV}}$ asociada a un precio observado $C^\star$ se define implícitamente como la solución de
\begin{equation}
C_{\mathrm{BS}}(S, K, T, r, \sigma_{\mathrm{IV}}) = C^\star,
\end{equation}
es decir, la $\sigma$ que habría que meter en Black-Scholes para reproducir ese precio. Comparar errores en precio puede ser muy engañoso porque un error de veinte céntimos no significa lo mismo en una opción que vale medio euro que en una que vale veinticinco, ni en una corta que en una larga, ni en una in-the-money que en una out-of-the-money. La volatilidad implícita actúa como unidad común que pone todas las opciones en una escala más homogénea. Aunque el modelo verdadero sea Bates o Heston, se toma el precio que genera y se busca qué volatilidad habría que meter en Black-Scholes para producir ese mismo precio, y luego se comparan volatilidades implícitas en lugar de precios. Esto es estándar en la práctica del mercado de opciones y es una decisión metodológica que deberíamos copiar tal cual.

\item \textbf{Generación de datos sobre un hipercubo con muestreo flexible.} Para la generación del dataset, se define el espacio de inputs como un hipercubo
\begin{equation}
\mathcal{X} = \prod_{i=1}^{d} [a_i, b_i],
\end{equation}
fijando mínimo y máximo para cada variable, se generan puntos aleatorios $\{x^{(j)}\}_{j=1}^{N} \sim \mathcal{X}$ dentro y para cada punto se evalúa el modelo verdadero $y^{(j)} = f(x^{(j)})$, formando así la muestra sintética $\mathcal{D} = \{(x^{(j)}, y^{(j)})\}_{j=1}^{N}$. La distribución de muestreo no tiene por qué ser uniforme: puede ser uniforme simple, basada en priors jerárquicos, o concentrarse en zonas donde la función es más no lineal o donde la red está cometiendo más error, lo cual conecta con active learning y es algo que nos gustaría explorar como elemento diferencial. Hay también un detalle técnico importante, y es que las redes suelen aproximar peor cerca de los bordes del dominio, así que los autores generan los datos en un rango ligeramente reducido, quitando un cinco por ciento por cada lado, para evitar que la evaluación quede dominada por errores de frontera.

\item \textbf{Lógica de entrenamiento distinta a la del ML clásico.} En machine learning normal se suele usar early stopping, dropout y otras técnicas de regularización para evitar overfitting. El entrenamiento estándar consiste en minimizar una pérdida empírica
\begin{equation}
\mathcal{L}(\theta) = \frac{1}{N} \sum_{j=1}^{N} \ell\!\left(\hat{f}_\theta(x^{(j)}),\, y^{(j)}\right),
\end{equation}
y en un surrogate no es necesario porque los datos son sintéticos y casi sin ruido, ya que el target viene del propio modelo. Si la red empieza a fallar en alguna zona siempre se pueden generar más datos en esos entornos. No estamos intentando aprender patrones inciertos a partir de datos limitados y ruidosos, sino aproximando una función determinista conocida pero cara de evaluar. Esa diferencia conceptual es la que justifica que los hiperparámetros y la lógica de entrenamiento sean distintos a los de un problema típico de ML.
\end{itemize}

\subsection{Cómo validar la precisión de los surrogates}

La validación que hacen en el paper también es interesante y está separada en dos niveles. En el primer nivel suponen que los parámetros verdaderos son conocidos y comparan precio del modelo Bates resuelto con FFT contra precio del surrogate. Para hacerlo de forma seria dividen las opciones en doce grupos según vencimiento y moneyness, con cuatro categorías de tiempo a vencimiento y tres niveles de moneyness ($4 \times 3 = 12$ bins), y dentro de cada grupo generan veinte mil opciones aleatorias. Esa estructura por bins permite ver si el surrogate funciona bien en distintas zonas del mercado, no solo en promedio. Sobre cada bin reportan el error absoluto medio y los percentiles de la distribución del error,
\begin{equation}
\mathrm{MAE}_{\mathrm{bin}} = \frac{1}{N_{\mathrm{bin}}} \sum_{j \in \mathrm{bin}} \left| \hat{f}_\theta(x^{(j)}) - f(x^{(j)}) \right|.
\end{equation}
Los resultados son muy buenos: el percentil noventa y cinco del error absoluto en volatilidad implícita y en Delta está por debajo de $0{,}0008$ para opciones con más de siete días a vencimiento. Los errores son algo mayores en opciones muy cortas, lo cual tiene sentido porque ahí el payoff está más cerca, hay más no linealidad y las sensibilidades son más bruscas, pero incluso ahí los errores siguen siendo pequeños en términos relativos.

En el segundo nivel de validación, una vez probado que la red aproxima bien los precios para parámetros dados, los autores van más allá y estudian si esos errores de aproximación se traducen en errores de estimación de parámetros cuando el surrogate se utiliza dentro de un procedimiento de calibración. La idea es que, en la práctica real, la estimación de parámetros del modelo se debe hacer a partir de los datos de mercado, y ese paso introduce una nueva fuente de error. Para aislar este efecto se realiza un experimento de simulación controlado, se fija un conjunto de parámetros verdaderos del modelo de Bates, se simula con el solver original una superficie de precios de opciones que ``actúan'' a modo de datos observados, y a continuación aplican el método generalizado de momentos (GMM) usando el surrogate como evaluador del modelo dentro del bucle de optimización para recuperar una estimación de esos parámetros. Una vez obtenida la estimación, miden el error de pricing fuera de muestra comparando los precios que salen al evaluar el solver original con los parámetros estimados frente a los precios que salen al evaluarlo con los parámetros verdaderos, lo que les permite cuantificar el coste económico atribuible exclusivamente al hecho de haber calibrado con la red en lugar de con el solver.

Hay además un detalle muy interesante que conviene destacar, algunos parámetros del modelo pueden estar débilmente identificados, en el sentido de que los precios de las opciones son muy poco sensibles a sus variaciones. En esos casos puede ocurrir que el error de estimación del parámetro en sí mismo parezca grande, porque pequeños movimientos en los precios observados se traducen en variaciones amplias del parámetro recuperado y sin embargo que el impacto económico de ese mismo error sobre los precios resultantes sea muy pequeño, precisamente porque la sensibilidad de los precios a ese parámetro también es baja.

\subsection{Aplicaciones prácticas}

Una parte del paper que se aleja de nuestro alcance pero que conviene entender es cómo el surrogate abre la puerta a aplicaciones empíricas que antes eran demasiado costosas. Como el surrogate permite reestimar Bates diariamente, los autores construyen un índice llamado TailRisk que mide la pérdida esperada bajo medida neutral al riesgo asociada a un salto negativo del índice durante la próxima semana. La idea es que las opciones contienen información sobre el riesgo de caídas extremas y los parámetros del modelo permiten cuantificar este riesgo. Después, comprueban si TailRisk predice eventos extremos reales del S\&P 500 mediante regresión logística y encuentran que el pseudo R cuadrado casi se duplica frente a una medida no paramétrica anterior basada en opciones.

Otra aplicación que hacen es estudiar la inestabilidad de los parámetros estimados. La intuición es que si los parámetros del modelo cambian mucho día a día, los market makers tienen más dificultad para cubrir sus inventarios, asumen más riesgo y exigen mayor compensación, y eso se acaba reflejando en bid-ask spreads más altos. Encuentran una relación positiva y estadísticamente significativa entre inestabilidad de parámetros y bid-ask spreads relativos en opciones SPX, incluso controlando por características del contrato, volumen anormal, vencimiento y moneyness. Esto sugiere que la inestabilidad del modelo no es solo un problema técnico, tiene consecuencias económicas reales sobre la liquidez del mercado.

Estos son dos ejemplos de aplicaciones prácticas que se desbloquean gracias a la técnica del surrogate y la reducción drástica que implica para el coste de las evaluaciones.

\section{El enfoque de nuestro trabajo}

Después de digerir todo esto, la decisión más importante es cómo enfocar nuestro trabajo para complementar y aportar valor a la bibliografía ya existente, dentro de nuestras limitaciones temporales y computacionales.

La estructura más coherente es trabajar en dos niveles. El primer nivel es Black-Scholes~\cite{blackscholes1973} como caso de validación. La fórmula cerrada para una call europea es
\begin{equation}
C_{\mathrm{BS}}(S, K, T, r, \sigma) = S\, \Phi(d_1) - K e^{-rT}\, \Phi(d_2),
\end{equation}
con
\begin{equation}
d_1 = \frac{\ln(S/K) + (r + \tfrac{1}{2}\sigma^2)T}{\sigma\sqrt{T}}, \qquad d_2 = d_1 - \sigma\sqrt{T},
\end{equation}
y la Delta analítica $\Delta_{\mathrm{BS}} = \Phi(d_1)$. Black-Scholes en la práctica no necesita un surrogate porque la fórmula cerrada ya es muy rápida, pero precisamente por eso es ideal como entorno de prueba. Conocemos la solución exacta, conocemos las Deltas analíticas, podemos comparar precio, volatilidad implícita y Greeks de forma directa, y podemos validar todo el pipeline antes de complicarlo. El segundo nivel es extender a Heston~\cite{heston1993}, cuya dinámica
\begin{align}
\frac{dS_t}{S_t} &= r\, dt + \sqrt{v_t}\, dW^S_t, \\
dv_t &= \kappa(\theta - v_t)\, dt + \xi \sqrt{v_t}\, dW^v_t, \qquad dW^S_t\, dW^v_t = \rho\, dt,
\end{align}
introduce volatilidad estocástica, es bastante más realista y su evaluación requiere métodos numéricos como Fourier semi-cerrado, así que la motivación de tener un surrogate empieza a cobrar sentido real. Es un paso intermedio entre Black-Scholes y Bates que conserva la esencia del paper sin meternos en el infierno de implementar saltos doble exponenciales.

La pregunta a responder es la siguiente:

\begin{quote}
Bajo qué condiciones una red neuronal profunda puede actuar como un surrogate preciso, rápido y diferenciable para funciones de pricing de opciones.
\end{quote}

Más allá de responder esta pregunta, hemos identificado cinco puntos en los que centrar nuestro proyecto de investigación. El primero es comparar de forma sistemática el error en precio frente al error en volatilidad implícita. La idea es plantearlo como una pregunta metodológica explícita: hasta qué punto puede un surrogate tener un error en precio aparentemente bajo pero un error relevante en implied volatility en ciertas zonas del dominio. Esto acerca el trabajo a la práctica real del mercado, donde la volatilidad implícita es la unidad natural de comparación.

El segundo elemento diferencial es comparar funciones de activación con foco en la calidad de los Greeks. Vamos a entrenar redes con la misma arquitectura cambiando solo la activación, probando ReLU, Softplus, SiLU/Swish y quizá tanh, y vamos a comparar el MAE de precio, el MAE de Delta y el MAE de volatilidad implícita en cada caso. La hipótesis, alineada con el paper, es que las activaciones suaves producen derivadas más estables aunque los errores de precio sean parecidos. Mostrar visualmente que dos redes con MAE de precio similar pueden tener Deltas muy distintas es una conclusión muy potente para un trabajo de métodos numéricos.

El tercer elemento es comparar muestreo uniforme frente a muestreo enfocado. Vamos a generar dos datasets, uno con muestreo uniforme en todo el dominio y otro con más densidad cerca del at-the-money y vencimientos cortos, que es donde sabemos por adelantado que la función de pricing tiene más curvatura. Después evaluaremos ambos surrogates con la misma división por bins y analizaremos si el muestreo enfocado reduce el error en regiones críticas. Esto conecta directamente con ideas clásicas de métodos numéricos como la elección adaptativa de puntos de malla y la concentración del esfuerzo computacional en zonas difíciles.

El cuarto elemento es la medición seria de eficiencia computacional. En Black-Scholes la ganancia probablemente sea pequeña porque la fórmula cerrada ya es rapidísima, pero el experimento sirve para validar el pipeline. Vamos a medir el tiempo necesario para valorar grandes volúmenes de opciones con el solver original frente al surrogate.

El quinto elemento, que solo abordaremos si vamos bien de tiempo después de cerrar los cuatro anteriores, es estudiar el efecto de la función de pérdida en la calidad del surrogate, con especial atención al enfoque de differential machine learning de Huge y Savine. Lo planteamos como opcional porque es el punto que más coste añadido tiene, requiere implementar el cálculo de Deltas durante la generación de datos, modificar la lógica de entrenamiento para que la pérdida combine error de precio y error de derivada, y diseñar comparaciones cuidadosas para que las conclusiones sean limpias.

La idea sería encadenar dos pruebas. La primera es comparar pérdidas clásicas, L1, L2 y Huber, manteniendo todo lo demás fijo, para ver cómo cambia la distribución del error por regiones:
\begin{align}
\mathcal{L}_{1}(\theta) &= \tfrac{1}{N} \textstyle\sum_j |\hat{f}_\theta(x^{(j)}) - y^{(j)}|, \\
\mathcal{L}_{2}(\theta) &= \tfrac{1}{N} \textstyle\sum_j (\hat{f}_\theta(x^{(j)}) - y^{(j)})^2, \\
\mathcal{L}_{\mathrm{Huber}}(\theta) &= \tfrac{1}{N} \textstyle\sum_j \mathrm{Huber}_\delta\!\left(\hat{f}_\theta(x^{(j)}) - y^{(j)}\right).
\end{align}
La hipótesis razonable, basada en lo que dice el paper de Chen sobre por qué ellos eligen $\ell_1$, es que $\ell_2$ dará mejor MAE en el centro del dominio pero peor en los bordes, y al revés con $\ell_1$, mientras que Huber debería comportarse como un compromiso. La segunda prueba, mucho más interesante, es probar la hipótesis de Huge y Savine aplicada a nuestro ámbito: en lugar de entrenar la red solo con precios, modificamos la pérdida para que también penalice la diferencia entre la Delta verdadera y la Delta que sale por autograd de la red,
\begin{equation}
\mathcal{L}_{\mathrm{diff}}(\theta) = \alpha \cdot \underbrace{\tfrac{1}{N} \textstyle\sum_j \left| \hat{f}_\theta(x^{(j)}) - y^{(j)} \right|}_{\text{error de precio}} + \beta \cdot \underbrace{\tfrac{1}{N} \textstyle\sum_j \left\| \nabla_x \hat{f}_\theta(x^{(j)}) - \nabla_x f(x^{(j)}) \right\|}_{\text{error de derivadas}}.
\end{equation}
De este modo, la pérdida queda como una combinación ponderada del error de precio y del error de la derivada y la red aprende simultáneamente a acertar el nivel y la pendiente de la función.

Este planteamiento tiene una motivación matemática sólida desde el punto de vista de aproximación de funciones, análoga a la que justifica la interpolación de Hermite frente a la clásica: disponer del valor y de la derivada de una función en cada punto muestreado permite reconstruirla con muchos menos puntos. La hipótesis asociada es que el differential machine learning ofrece una ventaja sustancial frente al entrenamiento con pérdidas estándar en regímenes de datasets pequeños, ventaja que se atenúa conforme crece el tamaño de la muestra hasta que las distintas configuraciones convergen, lo que permitiría cuantificar bajo qué condiciones compensa el coste adicional de calcular gradientes verdaderos durante la generación del dataset.

\section{Conclusiones}

Desde el primer día en nuestro grupo hemos decidido mantener un enfoque muy marcado en realizar una investigación profunda de la bibliografía existente y en comprender desde primeros principios el problema de fondo que aspiramos a resolver, convencidos de que esa base es la que permite tomar después decisiones técnicas informadas y bien fundamentadas. Conviene destacar esta elección dado que, en los tiempos que vivimos, la irrupción de la inteligencia artificial ha provocado que ejecuciones técnicas anteriormente muy tediosas y de bajo valor añadido se vuelvan prácticamente triviales, trasladando el grueso de la labor humana al plano de la comprensión profunda de los problemas subyacentes y de la planificación minuciosa del alcance, los métodos y las herramientas necesarias en cada proyecto. A lo largo de las próximas semanas pasaremos a ejecutar el trabajo apoyándonos en la palanca que nos ofrecen las herramientas hoy disponibles, con la confianza de que el tiempo invertido en esta fase preparatoria se traducirá en una implementación más sólida y mejor justificada.

\begin{thebibliography}{99}
\footnotesize
\setlength{\itemsep}{0.15em}
\setlength{\parskip}{0pt}

\bibitem{blackscholes1973}
Black, F. y Scholes, M. (1973).
\textit{The Pricing of Options and Corporate Liabilities}.
Journal of Political Economy, \textbf{81}(3), 637--654.

\bibitem{heston1993}
Heston, S.~L. (1993).
\textit{A Closed-Form Solution for Options with Stochastic Volatility with Applications to Bond and Currency Options}.
The Review of Financial Studies, \textbf{6}(2), 327--343.

\bibitem{bates1996}
Bates, D.~S. (1996).
\textit{Jumps and Stochastic Volatility: Exchange Rate Processes Implicit in Deutsche Mark Options}.
The Review of Financial Studies, \textbf{9}(1), 69--107.

\bibitem{hutchinson1994}
Hutchinson, J.~M., Lo, A.~W. y Poggio, T. (1994).
\textit{A Nonparametric Approach to Pricing and Hedging Derivative Securities via Learning Networks}.
The Journal of Finance, \textbf{49}(3), 851--889.

\bibitem{sirignano2018}
Sirignano, J. y Spiliopoulos, K. (2018).
\textit{DGM: A Deep Learning Algorithm for Solving Partial Differential Equations}.
Journal of Computational Physics, \textbf{375}, 1339--1364.

\bibitem{han2018}
Han, J., Jentzen, A. y E, W. (2018).
\textit{Solving High-Dimensional Partial Differential Equations Using Deep Learning}.
Proceedings of the National Academy of Sciences, \textbf{115}(34), 8505--8510.

\bibitem{ferguson2018}
Ferguson, R. y Green, A. (2018).
\textit{Deeply Learning Derivatives}.
arXiv preprint arXiv:1809.02233.

\bibitem{becker2019}
Becker, S., Cheridito, P. y Jentzen, A. (2019).
\textit{Deep Optimal Stopping}.
Journal of Machine Learning Research, \textbf{20}(74), 1--25.

\bibitem{buehler2019}
Buehler, H., Gonon, L., Teichmann, J. y Wood, B. (2019).
\textit{Deep Hedging}.
Quantitative Finance, \textbf{19}(8), 1271--1291.

\bibitem{huge2020}
Huge, B. y Savine, A. (2020).
\textit{Differential Machine Learning}.
arXiv preprint arXiv:2005.02347.

\bibitem{ruf2020}
Ruf, J. y Wang, W. (2020).
\textit{Neural Networks for Option Pricing and Hedging: A Literature Review}.
Journal of Computational Finance, \textbf{24}(1), 1--46.

\bibitem{dlcalib2023}
Sridi, A. y Bilokon, P. (2023).
\textit{Deep Learning Calibration of Stochastic Volatility Models via Differential Machine Learning}.
Working paper.

\bibitem{chen2025}
Chen, H., Didisheim, A. y Scheidegger, S. (2025).
\textit{Deep Surrogates for Finance: With an Application to Option Pricing}.
Working paper.

\end{thebibliography}

\end{document}
